\section{Introdução}
Um sistema de realimentação pode estabilizar uma planta com um polo no
semiplano direito. O mesmo laço pode perder estabilidade quando a comunicação
introduz atraso suficiente para alterar a fase relativa entre entrada e
saída. Esse efeito aparece em sistemas de controle em rede, nos quais o
tempo de transporte, a amostragem e a variação temporal da comunicação
afetam a malha \cite{zhang2001}. A pergunta investigada é em que condições
quatro métodos clássicos concordam sobre estabilidade e que erros emergem
quando a análise exige amostragem de frequências ou operações com
polinômios de ordem elevada.

A comparação usa uma implementação que mantém a análise de controle no
núcleo e isola as referências computacionais na validação. Os experimentos
avaliam o efeito de um polo instável, de um zero de fase não mínima e do
atraso puro. A ferramenta expõe os valores usados no diagnóstico e oferece
uma interface para variar os parâmetros e reproduzir os diagramas.

\section{Trabalhos relacionados}
Nyquist formulou a análise de regeneração em função da resposta de malha
\cite{nyquist1932}. Bode sistematizou o desenho de amplificadores realimentados
por magnitude e fase \cite{bode1945}. A contagem algébrica de Routh
\cite{routh1877} e o método geométrico de Evans \cite{evans1950} abordam a
mesma questão a partir do polinômio característico. Ogata e Nise organizam
esses métodos em uma sequência para sistemas contínuos
\cite{ogata2010,nise2017}. O capítulo de análise em frequência de Åström e
Murray relaciona Nyquist, margens e sensibilidade \cite{astrom2008}.
Para o atraso, a aproximação racional de Padé dá um caminho para métodos
polinomiais \cite{baker1996}. Em ordem elevada, a sensibilidade das raízes
a perturbações nos coeficientes pode dominar o erro observado
\cite{wilkinson1959}. Esses resultados orientam a escolha das referências
e das métricas deste estudo.

\section{Fundamentação teórica}
Considere a realimentação unitária negativa de uma planta contínua própria,
com coeficientes reais, ganho escalar $K$ e atraso $\tau\ge 0$.
\begin{equation}
L(s)=K\frac{B(s)}{A(s)}e^{-s\tau},\qquad
T(s)=\frac{L(s)}{1+L(s)} .
\label{eq:loop}
\end{equation}
Na parte racional, os polos e zeros são as raízes de $A$ e $B$,
respectivamente. Com $\tau=0$, o polinômio característico é
$C(s)=A(s)+KB(s)$. A análise preserva a ordem de $A$ e os fatores comuns,
para que um modo interno cancelado na análise entrada--saída permaneça
no modelo algébrico. Para estabilidade assintótica, os polos de malha
fechada devem ficar no semiplano esquerdo.

\subsection{Mapeamento e princípio do argumento}
O contorno segue o eixo de $-\imagunit R$ a $+\imagunit R$ e volta pelo arco à direita.
Indentam-se polos sobre o eixo com semicírculos que os excluem da
região contada. Para $F=1+L$, a variação de argumento no contorno
horário é $-2\pi(Z-P)$, onde $P$ e $Z$ contam polos e zeros no interior.
O mapeamento por $L$ escreve a mesma contagem em torno de $-1$.
\begin{equation}
\begin{aligned}
Z &= N+P,\\
N &= -\frac{1}{2\pi}\sum_{k}\arg
 \left(\overline{u_k}\,u_{k+1}\right),\\
u_k &= 1+L(s_k).
\end{aligned}
\label{eq:nyquist}
\end{equation}
A fórmula discreta fecha o polígono e preserva a convenção de que
envoltórios horários contribuam positivamente. Ela aproxima a
variação contínua quando o mapeamento entre amostras é resolvido
\cite{astrom2008}.

\subsection{Routh--Hurwitz, lugar das raízes e Bode}
A tabela de Routh dispõe os coeficientes alternados nas duas primeiras
linhas. Se $r_{i,j}$ é a entrada da linha $i$ e coluna $j$,
\[
r_{i,j}=r_{i-2,j+1}
-\frac{r_{i-2,0}}{r_{i-1,0}}\,r_{i-1,j+1}.
\]
As mudanças de sinal na primeira coluna contam raízes no semiplano
direito. Um pivô nulo exige tratar o limite de um valor positivo.
Uma linha inteira nula exige a derivada do polinômio auxiliar
\cite{routh1877,nise2017}. O lugar das raízes mostra o conjunto das
raízes de $A+KB$ para ganhos sucessivos e permite observar
amortecimento, frequência natural e cruzamentos \cite{evans1950}.

O diagrama de Bode mostra $20\log_{10}|L(\imagunit\omega)|$ e a fase
desenrolada. Em um cruzamento de fase $\omega_p$ no semieixo real
negativo, a margem de ganho é $1/|L(\imagunit\omega_p)|$. Em um cruzamento
de módulo unitário $\omega_g$, a margem de fase é o afastamento
angular de $-180^{\circ}$ na convenção local. Um ganho com margem
inferior a um pode indicar uma fronteira ao reduzir $K$. Quando
existem vários cruzamentos, cada um descreve uma fronteira candidata.
A classificação global usa a contagem de Nyquist ou os polos, conforme
o modelo \cite{bode1945,astrom2008}.

\subsection{Atraso e aproximação racional}
A exponencial acrescenta $-\omega\tau$ rad à fase e preserva o módulo
no eixo. A fronteira pode ocorrer em $\omega>0$ quando
$|L_0(\imagunit\omega)|=1$ e
\begin{equation}
\arg L_0(\imagunit\omega)-\omega\tau=(2q+1)\pi,
\qquad q\in\mathbb{Z}.
\label{eq:delay}
\end{equation}
Aqui $L_0$ é a função sem atraso. Para aproximar o atraso no plano
$s$, o Padé diagonal de ordem $m$ usa $Q_m(-s\tau)/Q_m(s\tau)$.
Se $c_k$ são os coeficientes crescentes de $Q_m$, então
$c_0=1$ e
\[
c_{k+1}=c_k\frac{m-k}{(k+1)(2m-k)}.
\]
A série coincide com a da exponencial até a potência $2m$.
O módulo dessa aproximação no eixo é unitário. A fase e os polos
fechados dependem da ordem e da banda de avaliação \cite{baker1996}.
